\section {Summary}
\label{sec:summary}

In this technical note, we have presented the results of the $\nueb$ appearance analysis using the 3-flavour $\nu$/$\bar{\nu}$ joint analysis framework performed by the VALOR group, on the combined \RunPeriod data set, for an exposure of \TotalPotNuRunOneNineC\ POT in neutrino mode and \TotalPotNuBarRunOneNineC\ POT in antineutrino mode. The goal of the analysis is to search for $\nueb$ appearance in a $\numub$ beam.

This analysis predicts $\NskOscExpectationNumuRunOneNineC \pm \NskOscExpectationNumuErrorRunOneNineC$ (syst) for neutrino $\mu$-like, $\NskOscExpectationNumuBarRunOneNineC \pm \NskOscExpectationNumuBarErrorRunOneNineC$ (syst) for antineutrino $\mu$-like, $\NskOscExpectationNueRunOneNineC \pm \NskOscExpectationNueErrorRunOneNineC$ (syst) for neutrino e-like, $\NskOscExpectationNueBarRunOneNineC \pm \NskOscExpectationNueBarErrorRunOneNineC$ (syst) for antineutrino e-like and $\NskOscExpectationNueCConepiRunOneNineC \pm \NskOscExpectationNueCConepiErrorRunOneNineC$ (syst) for neutrino $\nu_{e}\textrm{CC}1\pi^{+}$-like events in Super K and observes respectively \NskObservedNumuRunOneNineC\ for neutrino $\mu$-like, \NskObservedNumuBarRunOneNineC\ for antineutrino $\mu$-like, \NskObservedNueRunOneNineC\ for neutrino e-like, \NskObservedNueBarRunOneNineC\ for antineutrino e-like and \NskObservedNueCConepiRunOneNineC\ for neutrino $\nu_{e}\textrm{CC}1\pi^{+}$-like events.

Both a rate only and a rate+shape analysis are performed. In both analyses, fits were performed in the joint fit framework, including all samples as data to constraint either background and the hypotheses tested. The number of observed events in the antineutrino mode e-like spectrum is \NskObservedNueBarRunOneNineC\ and the measured rate+shape test statistic, $\Delta \chi^2 = \chi^2\left( \beta = 0 \right) - \chi^2\left( \beta = 1 \right)$ was found to be \DeltachisqData. The significance of the result is shown as a p-value. Taking no $\nueb$ appearance as our null hypothesis we find a p-value of \pvalRateOnlyBetaZeroData\ for the rate-only analysis and a p-value of \pvalRateShapeBetaZeroData\ for the rate+shape analysis. Taking $\nueb$ appearance according to the best knowledge of the PMNS matrix as our null hypothesis we find a p-value of \pvalRateOnlyBetaOneData\ for the rate-only analysis and a p-value of \pvalRateShapeBetaOneData\ for the rate+shape analysis. The Bayes factor, given by the rate+shape analysis was found to be \BayesFactorData\, indicating a poor preference for no $\nueb$ appearance.